\documentclass[UTF8,zihao=-4]{ctexart}
\usepackage[a4paper,margin=2.2cm]{geometry}
\usepackage{xcolor}
\usepackage{hyperref}
\usepackage{enumitem}
\usepackage{fontspec}
\usepackage{amsmath}
\IfFontExistsTF{Microsoft YaHei}{\setCJKmainfont{Microsoft YaHei}}{\setCJKmainfont{SimSun}[BoldFont=SimHei]}
\IfFontExistsTF{Microsoft YaHei UI}{\setCJKsansfont{Microsoft YaHei UI}}{}
\IfFontExistsTF{Consolas}{\setmonofont{Consolas}}{}
\definecolor{linkblue}{RGB}{30,80,145}
\hypersetup{colorlinks=true,linkcolor=linkblue,urlcolor=linkblue,pdfborder={0 0 0}}
\setlength{\parindent}{2em}
\setlength{\parskip}{0.25em}
\linespread{1.16}
\setlist{itemsep=0.2em,topsep=0.35em,leftmargin=2.5em}
\pagestyle{plain}

\title{11\_贪心法}

\author{}

\date{}

\begin{document}

\maketitle

\section*{11 贪心法}

\subsection*{题目原文}

贪心法的思想和特点，用贪心法解决经典问题并能给出证明思路：换硬币、单源最短路径。

\subsection*{考查类型}

概念题、算法设计题和证明题。

\subsection*{贪心法思想}

贪心法在每一步都做当前看来最优的选择，并希望这些局部选择组合成全局最优解。

适用贪心法通常需要两个性质：

\begin{quote}
\small\ttfamily\raggedright
\relax 贪心选择性质\\
\relax 最优子结构\\
\end{quote}

贪心选择性质表示存在一个全局最优解包含当前的局部最优选择。

最优子结构表示做出一次选择后，剩余问题的最优解可以和当前选择组合成原问题最优解。

\subsection*{证明思路}

常用证明方法有两类。

交换论证：

从任意最优解出发，把它逐步替换成包含贪心选择的最优解，同时不降低解的质量。

领先性证明：

证明贪心算法在每一步都不差于任何最优解的对应前缀。

\subsection*{换硬币问题}

目标是在给定币值系统下，用尽量少的硬币凑出金额 \texttt{N}。

贪心策略：

\begin{quote}
\small\ttfamily\raggedright
\relax 每次选择不超过剩余金额的最大面值硬币\\
\end{quote}

重要点：

不是所有币值系统都能用贪心得到最优解。答题时必须结合题目给出的币值系统证明贪心选择性质。

证明结构：

\begin{enumerate}
\item 说明贪心第一步选取最大可用面值。
\item 证明存在一个最优解也包含这枚硬币。
\item 删除这枚硬币后，剩余金额构成同类子问题。
\item 对子问题重复使用同样论证。
\end{enumerate}

若题目给出标准币值系统，需要用该系统的面值关系证明“大面值替换多个小面值不会增加硬币数”。

\subsection*{单源最短路径}

单源最短路径的典型贪心算法是 Dijkstra 算法。它要求边权非负。

算法思想：

\begin{enumerate}
\item 维护源点到各点的当前最短估计 \texttt{dist[v]}。
\item 每次从未确定顶点中取 \texttt{dist} 最小的顶点 \texttt{u}。
\item 将 \texttt{u} 标记为已确定。
\item 用 \texttt{u} 的出边松弛相邻顶点。
\end{enumerate}

伪代码：

\begin{quote}
\small\ttfamily\raggedright
\relax Dijkstra(G,\ s):\\
\relax \ \ \ \ for\ each\ vertex\ v:\\
\relax \ \ \ \ \ \ \ \ dist[v]\ =\ ∞\\
\relax \ \ \ \ dist[s]\ =\ 0\\
\end{quote}
\[Q = all vertices\]
\begin{quote}
\small\ttfamily\raggedright
\mbox{}\\
\relax \ \ \ \ while\ Q\ is\ not\ empty:\\
\relax \ \ \ \ \ \ \ \ u\ =\ vertex\ in\ Q\ with\ minimum\ dist[u]\\
\relax \ \ \ \ \ \ \ \ remove\ u\ from\ Q\\
\relax \ \ \ \ \ \ \ \ for\ each\ edge\ (u,\ v):\\
\relax \ \ \ \ \ \ \ \ \ \ \ \ if\ dist[v]\ >\ dist[u]\ +\ w(u,\ v):\\
\relax \ \ \ \ \ \ \ \ \ \ \ \ \ \ \ \ dist[v]\ =\ dist[u]\ +\ w(u,\ v)\\
\end{quote}

\subsection*{Dijkstra 正确性证明}

关键命题：

每次取出的 \texttt{dist} 最小未确定顶点 \texttt{u}，其 \texttt{dist[u]} 已经等于源点到 \texttt{u} 的真实最短距离。

证明思路：

若存在一条更短路径从源点到 \texttt{u}，考虑这条路径上第一个未确定顶点 \texttt{y}，以及它前一个已确定顶点 \texttt{x}。由于 \texttt{x} 已经确定，边 \texttt{(x, y)} 已经被松弛，因此 \texttt{dist[y]} 不会大于该路径到 \texttt{y} 的长度。边权非负，所以到 \texttt{y} 的长度不大于到 \texttt{u} 的长度。这会与 \texttt{u} 是未确定顶点中 \texttt{dist} 最小者矛盾。

因此贪心选择正确。

\subsection*{例题}

\subsubsection*{例题 1}

题目：币值为：

\begin{quote}
\small\ttfamily\raggedright
\relax 1,\ 2,\ 4,\ 8,\ 16\\
\end{quote}

用贪心法凑出金额 \texttt{23}，并说明最优性。

解答：

贪心选择：

\[23 = 16 + 4 + 2 + 1\]

共用 4 枚硬币。

这些币值都是 2 的幂。金额 \texttt{23} 的二进制表示为：

\[23 = 16 + 4 + 2 + 1\]

每个二进制位对应一枚硬币。使用更小面值替代其中任意一枚硬币，都不会减少硬币数。因此该贪心解是最优解。

\subsubsection*{例题 2}

题目：对下面非负权图从源点 \texttt{s} 运行 Dijkstra 算法：

\[s->a: 1\]
\[s->b: 4\]
\[a->b: 2\]
\[a->c: 5\]
\[b->c: 1\]
\[b->d: 7\]
\[c->d: 3\]

求从 \texttt{s} 到各点的最短距离。

解答：

初始化：

\begin{quote}
\small\ttfamily\raggedright
\relax dist[s]\ =\ 0\\
\relax dist[a]\ =\ ∞\\
\relax dist[b]\ =\ ∞\\
\relax dist[c]\ =\ ∞\\
\relax dist[d]\ =\ ∞\\
\end{quote}

从 \texttt{s} 松弛：

\begin{quote}
\small\ttfamily\raggedright
\relax dist[a]\ =\ 1\\
\relax dist[b]\ =\ 4\\
\end{quote}

取出 \texttt{a}，松弛：

\begin{quote}
\small\ttfamily\raggedright
\relax dist[b]\ =\ min(4,\ 1\ +\ 2)\ =\ 3\\
\relax dist[c]\ =\ 1\ +\ 5\ =\ 6\\
\end{quote}

取出 \texttt{b}，松弛：

\begin{quote}
\small\ttfamily\raggedright
\relax dist[c]\ =\ min(6,\ 3\ +\ 1)\ =\ 4\\
\relax dist[d]\ =\ 3\ +\ 7\ =\ 10\\
\end{quote}

取出 $c$，松弛：

\begin{quote}
\small\ttfamily\raggedright
\relax dist[d]\ =\ min(10,\ 4\ +\ 3)\ =\ 7\\
\end{quote}

最终结果：

\begin{quote}
\small\ttfamily\raggedright
\relax dist[s]\ =\ 0\\
\relax dist[a]\ =\ 1\\
\relax dist[b]\ =\ 3\\
\relax dist[c]\ =\ 4\\
\relax dist[d]\ =\ 7\\
\end{quote}

\subsection*{常见误区}

贪心算法必须证明正确性，不能只写“每一步取最优”。

Dijkstra 算法不能直接用于含负权边的图。

换硬币问题不能默认所有币值系统都适合贪心。

\end{document}
